\section{The refined coefficient-three family on the same window}
\label{sec:coefficient-three-refined-window}

\paragraph{The previous union and the original completion are retained.}
Keep every datum, field, page, prescribed pair and coefficient ring
of AC-G176-window.  In particular
\[
 \begin{gathered}
 H_u=\mathbb Q(\zeta_{2^{h(u)+2}},\zeta_7,
             \sqrt{-q},\sqrt{-d},\omega,\sqrt[3]{269}),\qquad
 K=\mathbb Q(\omega),\\
 u^{2+8\varepsilon}\leq2^{h(u)}\leq2u^{2+8\varepsilon},
 \quad 0<\varepsilon<1/256,\quad T=\lfloor u^{399}\rfloor.
 \end{gathered}\tag{R3W.1}
\]
For $u\geq2$ one has $h(u)\geq2$.
Write $Z^{\rm odd}$ and $I_{399}^{\rm unifodd}$ for the exact
omitted set and remainder of (OPW.3)--(OPW.6).
All their coefficient primes other than three are left unchanged.
The coefficient-three completion $O_3$ is the one already chosen
in (FC.2); its residue field is denoted $k_3$.
No new choice of coefficient embedding or lattice is made below.

\paragraph{The enlarged set of pure coefficient-three labels.}
For an integer $h\geq2$ let $\mathscr D_h$ consist of
\[
 \begin{split}
 \mathscr D_h={}&\{\chi_{-4},\chi_8\}|_{G_K}\\
 &\sqcup\{\delta|_{G_K}:\operatorname{cond}\delta=16\}\\
 &\sqcup\{\delta|_{G_K}:\operatorname{cond}\delta=32,
                      \ \delta(-1)=-1,\ h\geq3\}\\
 &\sqcup\{\delta|_{G_K}:\operatorname{cond}\delta=2^c,
                                      \ 6\leq c\leq h+2\}.
 \end{split}\tag{R3W.2}
\]
Here and below a character of a stated conductor is primitive.
Define also, using reduction in this actual $O_3$,
\[
 \mathscr V_h(O_3)=
 \left\{\delta|_{G_K}:\begin{array}{l}
 h\geq3,\ \operatorname{cond}\delta=32,\ \delta(-1)=1,\\
 \beta=\overline{\delta(5)}\in k_3,\quad\beta^2+\beta-1=0
 \end{array}\right\},\qquad
 \mathscr E_h(O_3)=\mathscr D_h\sqcup\mathscr V_h(O_3).
 \tag{R3W.3}
\]
The defining polynomial selects exactly two even conductor-32
characters for each original completion when $h\geq3$.
It does not permit replacing that completion by a convenient one.
The membership condition is invariant under residue-field Frobenius.
All these labels have trivial conductor-seven, auxiliary quadratic
and cubic factors, and are nontrivial.

Put
\[
 \begin{aligned}
 B(h;M)&=\max\{0,2^{h+1}-2^{M-2}\},\qquad N_T=T-1,\\
 C_3(h)&=B(h;8),\\
 D_3(h)&=6+4\mathbf1_{h\geq3}+B(h;6),\\
 E_3(h)&=D_3(h)+2\mathbf1_{h\geq3}.
 \end{aligned}\tag{R3W.4}
\]
$E_3(h)$ is a character count.
Define the new omitted set by the ordinary set union
\[
 Z^{\rm ref3}(u)=Z^{\rm odd}(u)\ \cup\
       \{(3,\eta,e):\eta\in\mathscr E_{h(u)}(O_3),\ 1\leq e<T\},
 \tag{R3W.5}
\]
and let $I_{399}^{\rm ref3}$ omit precisely this set from (TC.1).
The old coefficient-three trivial label remains in $Z^{\rm odd}$.
Its original-coefficient source is AC-K11 and the transport
(EW.8)--(EW.9) of AC-G146, retained through G176.

\begin{theorem}[AC-G179-window: the refined coefficient-three union]
\label{AC-G179-window}
Every triple in (R3W.5) is an actual original-coefficient strict zero.
The old coefficient-three high-conductor set is contained in
$\mathscr D_h$, so all previous omissions are retained.
For positive integer $T$,
\[
 \begin{aligned}
 |Z^{\rm ref3}\setminus Z^{\rm odd}|
                    &=(E_3(h)-C_3(h))N_T,\\
 |Z^{\rm ref3}_3|&=(1+E_3(h))N_T,\\
 |Z^{\rm ref3}_{269}|&=(8+C(h))N_T-B_T+12A_T-R_T-S_T,\\
 |Z^{\rm ref3}_\ell|&=C_\ell(h)N_T\qquad(\ell\in\Pi(h)),\\
 |Z^{\rm ref3}|&=(9+C(h)+E_3(h)+P(h))N_T
                                  -B_T+12A_T-R_T-S_T.
 \end{aligned}\tag{R3W.6}
\]
Every symbol other than $D_3,E_3$ in this display retains exactly
its definition in (OPW.2)--(OPW.4).
One has $E_3(2)=6$ and
\[
 E_3(h)=2^{h+1}-4\quad(h\geq3),\qquad
 E_3(h)-C_3(h)=
 \begin{cases}6,&h=2,\\12,&h=3,\\28,&h=4,\\60,&h\geq5.
 \end{cases}\tag{R3W.7}
\]
In particular the difference is $60$ for every $h\geq6$.
For $h\geq3$ the proved coefficient-three set, including the
old trivial label, contains all but three pure two-power labels:
the omitted ones are $\chi_{-8}$ and the two even primitive
conductor-32 labels failing (R3W.3).
Those three strict groups are not asserted nonzero.
The enlarged set is proper for $h\geq2$, $T\geq2$.

For the normalized positive-index cost $Q_u$ of the entire new union,
\[
 \begin{aligned}
 I_{399}&=I_{399}^{\rm ref3}+Q_u,\\
 I_{399}^{\rm ref3}&\leq I_{399}^{\rm unifodd}\leq I_{399},\\
 0\leq Q_u&\leq8D_u^2T^2\log(8TD_u)
                +\tfrac12D_u\mathcal R_uT^2+2D_ur_uc_0T\\
 &=O_\varepsilon(u^{802+16\varepsilon}(\log u)^2).
 \end{aligned}\tag{R3W.8}
\]
These are nonnegative extended sums, with one full (FL.4) cost,
not an additional cost for each old or new subfamily.
The same common-window implication with positive relative slack,
all original coefficients and full local conditions, and actual
source rank $R.t$ holds.  No individual derived index is asserted zero.
\end{theorem}

\begin{proof}
\proofstep{1. Literal strict theorems on the actual original ring.}
Every character in (R3W.2)--(R3W.3) has conductor dividing
$2^{h+2}$ and values of order dividing $2^h$.
It therefore occurs in the actual host and has values already in
the prescribed completion of $\mathbb Z[\zeta_{3\cdot2^h}]$.
The coefficient is canonically
\[
 (O_3\otimes_{\mathbb Z_3}\mathbb Q_3/\mathbb Z_3)(\eta)(-e)
                         =(E_3^{\rm fld}/O_3)(\eta c_3^{-e}),
 \tag{R3W.9}
\]
where $E_3^{\rm fld}=\operatorname{Frac}(O_3)$.
For $\eta\in\mathscr D_h$, apply (G178.2) with this original ring,
$s=0$ and $\kappa=1$.
For $\eta\in\mathscr V_h(O_3)$, apply (G180.2) with the same ring
and the same parameters.  The reduction condition in (G180.1)
is exactly the one imposed in (R3W.3), at the chosen completion.
Both theorems permit arbitrary original coefficient ramification
and every positive Tate weight.
The certified real class-group input in G180 is retained; this
application does not deduce its two rows from G178's nonunit
annihilators alone.

The strict condition above three is full local zero, and away
three it is zero on inertia, exactly as in (FC.2).
Inflation from its prescribed ramification group into $G_K$
is injective on this unchanged module.
Outside that ramification set an inflated cocycle already kills
inertia; inside it the same full local and inertia conditions hold.
Thus every new actual strict group injects into the corresponding
zero group of G178 or G180.
All old strict zeros are retained from AC-G176-window.

\proofstep{2. Exact disjointness and counts at the chosen completion.}
Distinct two-power characters remain distinct on $G_K$, since
their rational ratio cannot be the conductor-three character
$\chi_{-3}$.  The parts of $\mathscr D_h$ are separated by
their conductors and parity, and $\mathscr V_h$ has the even
parity at conductor32 missing from $\mathscr D_h$.
There are two named quadratic labels, four primitive conductor16
labels and four odd primitive conductor32 labels.
The primitive conductor-$2^c$ count is $2^{c-2}$.
Summing for $6\leq c\leq h+2$ gives $B(h;6)$ and hence $D_3(h)$.

For the four even primitive conductor32 characters the value at5
is each of the four primitive eighth roots once.
The factorization
$X^4+1=(X^2+X-1)(X^2-X-1)$ in characteristic three has distinct
quadratic factors, each with two distinct roots.
Reduction preserves eighth-root order, so exactly two labels
satisfy (R3W.3) for every actual $O_3$.
This proves $E_3(h)$ without changing the embedding or choosing
different fields or windows for the selected labels.
The old $c\geq8$ family is a subset of the new $c\geq6$ family,
and every new label is nontrivial.
Consequently the coefficient-three difference is exactly
$(E_3-C_3)N_T$; no other coefficient row changes.
Adding the unchanged (OPW.5) counts proves (R3W.6).
Direct substitution gives (R3W.7).
The complete two-power character group has $2^{h+1}$ elements;
for $h\geq3$ subtracting $1+E_3(h)$ leaves precisely the three
stated unsupported labels.  This is a description of the proved
set and makes no claim about nonvanishing outside it.

\proofstep{3. One full-union error on the original normalization.}
Keep the original dual codifferent and normalizing rank $b_\ell$.
For every omitted triple, (EU.2) and strict zero give
\[
 \frac{\log\#k_{O_\ell}}{b_\ell}\max(0,i_{\ell,\eta,e})
 \leq\frac{|\log\#H^2(C_{\ell,\eta,e})
                    -\log\#\mathcal S_{\ell,\eta}(e)|}{b_\ell}.
 \tag{R3W.10}
\]
At each weight the complete union is one subsum of the
nonnegative all-prime, all-character (FL.4) sum.
Summing its bound once over $1\leq e<T$, with
$\sum e\leq T^2/2$, proves (R3W.8).
The unchanged bounds $D_u\leq288u^{2+8\varepsilon}$,
$r_u\leq12$ and
$\mathcal R_u=O_\varepsilon(u^{\varepsilon/5}(\log u)^2)$
give the displayed exponent and uniformity in the prescribed pair.
Nonnegative additivity proves the extended-sum decomposition
without subtracting infinities.

\proofstep{4. The identical common window.}
Fix the retained finite datum, page and $\varepsilon$.
Suppose the new remainder satisfies the same (TC.3) quantifiers:
$J\geq1$, $X_0$ and $s\in(0,1)$ are independent of
$X,q,d,u,\ell,\eta,e$, and every prescribed pair at $X\geq X_0$
has one $X^{1/J}\leq W$, $W^2\leq X$, with
$I_{399}^{\rm ref3}(H_u,u)\leq(1-s)\exp(u^\varepsilon/2)$
throughout $[W,W^2]$.
The error (R3W.8) is $o(\exp(u^\varepsilon/2))$ uniformly
in that pair.  Choose $U_0\geq2$ making it at most half the
slack, replace $X_0$ by $\max(X_0,U_0^J)$, and use exactly
the already supplied $W$.
The full original premise holds there with slack $s/2$.
Conversely every older premise implies the new one by nonnegativity.
Uniform little-$o$ and fixed finite simultaneous families use
the same argument and a maximum of their uniform error constants;
no independently chosen windows are intersected.
All original coefficients, full local conditions, actual $R.t$
and simultaneous field/weight/character quantifiers remain those
of AC-K35.  Its mixed input $\mathcal R_{\rm cub}$ is unchanged.
\end{proof}

\paragraph{Boundary.}
All other-prime rows and every previously unsupported mixed or
auxiliary label remain exactly as in G176.
The theorem gives no full numerical remainder estimate, no
individual index-zero assertion, and no mixed realization or H04.4
conclusion.  It adds no stronger AW.1 premise.
The finite-cover, Artin--Tate and Kummer dictionary is the already
audited SGA~1, Expos\'e~V, \S\S4,7 and Stacks~03PK scope.
No new motive, archive attribution or Lean execution is claimed.
